\documentclass[11pt]{article}
\usepackage[a4paper,margin=1in]{geometry}
\usepackage[T1]{fontenc}
\usepackage[utf8]{inputenc}
\usepackage{lmodern}
\usepackage{amsmath,amssymb,amsthm}
\usepackage{booktabs}
\usepackage{array}
\usepackage[numbers,sort&compress]{natbib}
\usepackage[hidelinks]{hyperref}

\newtheorem{proposition}{Proposition}
\newtheorem{theorem}{Theorem}
\newcommand{\Vact}{V_{\mathrm{act}}}
\newcommand{\Iact}{I_{\mathrm{act}}}
\newcommand{\SG}{S_{\mathrm{G}}}
\newcommand{\STE}{S_{\mathrm{TE}}}
\newcommand{\SDI}{S_{\mathrm{DI}}}
\newcommand{\DT}{D_{\mathrm{T}}^{\ast}}
\newcommand{\DVKW}{\Delta V_{\mathrm{KW}}}
\newcommand{\SKW}{S_{\mathrm{KW}}}
\newcommand{\JS}{\mathrm{JS}}
\newcommand{\Oscr}{\Omega_{\mathrm{scr}}}
\newcommand{\PhiTC}{\Phi_{\mathrm{TC}}}
\newcommand{\Von}{V_{\mathrm{on}}}
\newcommand{\Voff}{V_{\mathrm{off}}}
\newcommand{\tr}{\mathrm{tr}}
\newcommand{\TC}{\mathrm{TC}}

% DRAFT 1, 2026-11-26. Every number comes from docs/B.md, docs/C.md and the files in results/.
% Status: not compiled on this PC (no TeX). Compile on Overleaf with pdfLaTeX.

\title{Semantic information in linear Gaussian systems:\\
closed forms, bounds and negative controls}
\author{Alexander Beiro\\ \small Independent researcher}
\date{}

\begin{document}
\maketitle

\begin{abstract}
Semantic information, in the sense of Kolchinsky and Wolpert, is the part of the information about its environment
that a system needs to keep its viability. Its computation needs a search over interventions. Earlier work names this
cost as the main obstacle. We compute it for linear Gaussian systems. For the stored form,
Gaussian pre-garblings of the environment give an exact formula for the viability as a function of one positive
semidefinite matrix $P$, with $0\preceq P\preceq I$. For an agent with one coordinate, the stored semantic information
has a closed form. It is the information of the agent about the one direction of the environment that drives it. For
more coordinates, noisy garblings beat deterministic coarse-grainings in 40 of 48 test systems, and the minimum needs a
search over one matrix. Non-Gaussian partitions beat the Gaussian value in 2 of 27 scalar systems, so the Gaussian
value is an upper bound. For the observed form, we prove that one time step never changes the viability and that the
directed information is at most the transfer entropy. Numerically, the transfer entropy of the agent is at most the sensor
information in the measure of Tanaka et al.\ at all 1508 tested sensors. So the observed semantic information is at
most the smallest such sensor information, and about six times smaller in the median. Then we apply the formulas and Causal Leverage Density (CLD) to two negative controls: a homeostat whose
actuator does not read its sensor, and passive bistable units. No score gives zero. For every organism without feedback, the value of
information of Kolchinsky and Wolpert is negative (we prove this), and the stored semantic information is the full
mutual information. CLD rates passive units above driven oscillators of the same life (stratified AUC 0.711 [0.676,
0.748]). These scores measure correlation or memory. A test of use needs a term that sees the drive. We fixed
all rules before the runs.
\end{abstract}

\section{Introduction}
A living system holds information about its environment, but not all of it matters. Kolchinsky and Wolpert proposed
to call \emph{semantic} the part of the information that the system needs to keep its viability
\cite{KolchinskyWolpert18}. They scramble the correlations between a system $X$ and its environment $Y$, and they
measure how the viability changes. The smallest amount of information that keeps the actual viability is the semantic
information. The idea is now part of a programme on information as a feature of life \cite{Bartlett25}, with
applications to foraging agents \cite{Sowinski23} and to a feedback-controlled Kuramoto model
\cite{SowinskiFrankGhoshal24}.

The computation is the obstacle. The definition needs a search over all coarse-grainings of the environment, and two
groups name its cost as an obstacle \cite{Sowinski23,SowinskiFrankGhoshal24}. One of them puts the limit at several
hundred states \cite{SowinskiFrankGhoshal24}.
Both groups therefore replaced the search with a one-parameter family of interventions (sensor noise). In a search of
the 181 papers that cite the main sources, we found no computation of semantic information for continuous or Gaussian
variables.

This paper gives that computation for linear Gaussian systems, and it tests the resulting scores on negative controls.
It has four parts.
\begin{enumerate}
\item \textbf{Stored semantic information} (Sec.~\ref{sec:static}). A Gaussian pre-garbling of the environment is one
matrix $P$ with $0\preceq P\preceq I$. The viability is an exact function of $P$. For an agent with one coordinate the
semantic information has a closed form. For more coordinates it needs a search over $P$, which takes seconds.
\item \textbf{Observed semantic information} (Sec.~\ref{sec:dynamic}). With interventions at each time step, we prove
two general facts and compare the result with the minimal directed information of linear quadratic Gaussian control
\cite{TanakaEsfahaniMitter18}.
\item \textbf{Application} (Sec.~\ref{sec:apply}). We apply the formulas to a linear homeostat and to the structures of
random linear worlds from a companion study \cite{Beiro26}.
\item \textbf{Negative controls} (Sec.~\ref{sec:controls}). A score of the use of information must give zero, or a
low value, to a system that holds information but does not use it. We test three scores on two such systems: the value of
information and the semantic information of Kolchinsky and Wolpert, and the Causal Leverage Density (CLD) of Bartlett
and Sowinski et al.\ \cite{Bartlett24,SowinskiEtAl25}.
\end{enumerate}

\paragraph{Pre-registration.} Each step had written rules, pass conditions and expectations before its runs. The rules
are in the step documents of the repository (\texttt{docs/B.md}, \texttt{docs/C.md}), each in a commit before the
commit of the results. We report every verdict as the rule gives it, also when the rule was badly written
(Sec.~\ref{sec:b1}). Values in brackets are 95\% bootstrap intervals \cite{Efron79}. Information is in nats, except
CLD, which is in bits as in its definition.

\section{Setting}
\label{sec:setting}
\paragraph{Dynamics.} The joint state $z=(x,y)$ has an agent part $x\in\mathbb{R}^{d_x}$ and an environment part
$y\in\mathbb{R}^{d_y}$. It follows the Ornstein--Uhlenbeck process
\begin{equation}
  \mathrm{d}z = -Bz\,\mathrm{d}t + \sqrt{2D}\,\mathrm{d}W
\end{equation}
\cite{Gardiner04}, and it starts from its stationary law $z_0\sim\mathcal{N}(0,S_0)$. The propagator is
$\Phi=e^{-B\tau}$, and the noise covariance $Q_\tau=\int_0^\tau e^{-Bs}2De^{-B^\top s}\mathrm{d}s$ comes from the block
exponential of Van Loan \cite{VanLoan78}.

\paragraph{Viability.} As in Ref.~\cite{KolchinskyWolpert18}, Eq.~(2), the viability is the negative entropy of the agent
at a time $\tau$:
\begin{equation}
  V = -H(X_\tau) = -\tfrac12\bigl(d_x\log 2\pi e + \log\det\Sigma_{X_\tau}\bigr).
\end{equation}
This choice rewards a narrow distribution of the agent. It is not the exit-time viability of the companion study
\cite{Beiro26}. We return to this difference in Sec.~\ref{sec:apply}.

\paragraph{Interventions.} Ref.~\cite{KolchinskyWolpert18} replaces the channel $p(x_0\mid y_0)$ with
$p(x_0\mid\varphi(y_0))$, for a coarse-graining $\varphi$ of the environment. The marginals of $X_0$ and $Y_0$ do not
change, and the dynamics after time 0 do not change. The full scrambling $p(x_0)p(y_0)$ is the coarsest case. The
paper names the general class \emph{pre-garbling}, but it computes only deterministic $\varphi$. We use the Gaussian
pre-garblings $Z=KY+N$, $N\sim\mathcal{N}(0,\Sigma_N)$, with $x_0$ drawn from the actual law $p(x_0\mid z)$.

\section{Stored semantic information}
\label{sec:static}

\subsection{The viability as a function of one matrix}
\label{sec:b1}
Write $C=\Sigma_X^{-1/2}\Sigma_{XY}\Sigma_Y^{-1/2}$ for the whitened cross-covariance at time 0.

\begin{proposition}[Gaussian pre-garbling]\label{prop:P}
Let $P=K_w^\top(K_wK_w^\top+\Sigma_N)^{-1}K_w$, with $K_w=K\Sigma_Y^{1/2}$. Then $0\preceq P\preceq I$, and every such
$P$ comes from some $(K,\Sigma_N)$. The intervention keeps the marginals of $X_0$ and $Y_0$ and changes only the
cross-covariance:
\begin{equation}
  \Sigma_{XY}(P)=\Sigma_{XY}\Sigma_Y^{-1/2}P\,\Sigma_Y^{1/2}, \qquad
  I_P(X_0;Y_0) = -\tfrac12\log\det\bigl(I-CP^2C^\top\bigr).
\end{equation}
The covariance of the agent at $\tau$ is affine in $P$:
\begin{equation}
  \Sigma_{X_\tau}(P) = A_0 + M(P) + M(P)^\top, \qquad M(P)=\Phi_{XX}\,\Sigma_{XY}(P)\,\Phi_{XY}^\top ,
  \label{eq:affine}
\end{equation}
where $A_0$ does not depend on $P$. So $V(P)$ is convex in $P$, because $\log\det$ is concave.
\end{proposition}

The proof is a direct computation of the Gaussian conditional law. $P=I$ is the actual system, $P=0$ is the full
scrambling, and the orthogonal projections are the deterministic coarse-grainings $\varphi(y)=U^\top\Sigma_Y^{-1/2}y$.
The information $I_P$ equals $I(X_0;Z)$ only when $P$ is a projection. Two consequences follow at once. First, if
the environment does not drive the agent up to time $\tau$ ($\Phi_{XY}=0$), then $V$ does not depend on $P$. The agent
then has no stored semantic information, whatever its mutual information with the environment. Second, the set of $P$
is compact and convex. So the search is over one $d_y\times d_y$ matrix, not over partitions of a state space.

\paragraph{Check by simulation.} We checked Proposition~\ref{prop:P} against an independent simulation: the start
sampled through the garbling itself, and Euler--Maruyama dynamics with $\mathrm{d}t=5\cdot10^{-4}$. There were 4 systems
($d_x=2$, $d_y=3$), 2 values of $\tau$, 4 interventions (actual, full scrambling, a random rank-1 projection and a
random garbling) and $10^5$ trajectories in 20 batches. The rule asked for at least 45 of 48 checks within 3 standard
errors (SE) and none beyond 5 SE. The verdict is \emph{fail}: 44 of 48. All 32 checks of $V$ pass (largest
$|z|=2.81$), and all 12 checks of $I_P$ with a nonzero true value pass. The 4 failures are $I_P$ at the full
scrambling, where the true value is 0. There, the plug-in estimate of a Gaussian mutual information has a known upward
bias of about $d_xd_y/2n=0.0006$ nats, which is 7 to 9 SE. The rule did not allow for this bias. After a subtraction of
this bias (a check made after the run), all 16 checks of $I_P$ are within 2.42 SE.

\subsection{The minimum}
The stored semantic information is the smallest information that keeps the actual viability
\cite{KolchinskyWolpert18}, Eqs.~(9)--(10). Over Gaussian garblings,
\begin{equation}
  \SG = \min\bigl\{\,I_P(X_0;Y_0)\;:\;0\preceq P\preceq I,\; V(P)=V(I)\,\bigr\},
\end{equation}
and the semantic efficiency is $\SG/\Iact$, with $\Iact=I_{I}(X_0;Y_0)$.

Let $W=\Phi_{XY}\Sigma_Y^{1/2}$. Its rows are the directions of the whitened environment that drive the agent up to
time $\tau$. The projector $\Pi_W$ on the row space of $W$ keeps $\Sigma_{X_\tau}$ exactly, because
$M(\Pi_W)=M(I)$. So the \emph{drive information} $I_C=-\tfrac12\log\det(I-C\Pi_WC^\top)$ is always an upper bound,
$\SG\le I_C$.

\begin{theorem}[one agent coordinate]\label{thm:dx1}
Let $d_x=1$, let $c=C^\top\in\mathbb{R}^{d_y}$, and let $w=W^\top$ with $c^\top w\neq0$. Then
\begin{equation}
  \SG = I_C = -\tfrac12\log\bigl(1-(c^\top\hat w)^2\bigr), \qquad \hat w = w/\lVert w\rVert .
\end{equation}
\end{theorem}
\begin{proof}
With $d_x=1$, Eq.~(\ref{eq:affine}) makes $\Sigma_{X_\tau}$ a scalar affine in $c^\top Pw$. If
$\Phi_{XX}\neq0$, then $V(P)=V(I)$ means $c^\top Pw=c^\top w$. Also $I_P=-\tfrac12\log(1-\lVert Pc\rVert^2)$ grows with $\lVert Pc\rVert$. The Cauchy--Schwarz
inequality gives $c^\top w=(Pc)^\top w\le\lVert Pc\rVert\,\lVert w\rVert$, so $\lVert Pc\rVert^2\ge(c^\top\hat w)^2$.
The projection $P=\hat w\hat w^\top$ reaches this bound and keeps $c^\top Pw=c^\top w$.
\end{proof}

So for an agent with one coordinate, the stored semantic information is the information of the agent about the one
direction of the environment that drives it. For $d_x\ge2$, the condition $V(P)=V(I)$ fixes only $\det\Sigma_{X_\tau}$,
not the whole matrix, and the proof does not carry over.

\paragraph{Tests.} We used SLSQP with 32 random starts, plus $P=I$ and $P=\Pi_W$, at $\tau=1$
(Table~\ref{tab:b2}). A sanity rule checked the optimizer against Theorem~\ref{thm:dx1} on 20 systems with $d_x=1$.
For non-Gaussian interventions, we searched three families of partitions of a scalar environment by brute force. The
families are one threshold, the two cells $|y|<t$ and $|y|\ge t$, and two thresholds. A bisection refines each sign
change of $V-\Vact$. With $d_y=1$, the only Gaussian garbling that keeps $V$ is $P=1$, so there $\SG=\Iact$.

\begin{table}[t]
\centering
\caption{Stored semantic information: pre-registered rules. Systems with $|\Delta V|<10^{-3}$ nats are skipped.}
\label{tab:b2}
\begin{tabular}{p{6.2cm}p{5.2cm}p{2.6cm}}
\toprule
Rule & Result & As expected? \\
\midrule
Sanity: $\SG=I_C$ for $d_x=1$ (19 systems) & pass, largest gap $1.8\cdot10^{-10}$ nats & yes \\
H1: projections reach the minimum ($d_x=2$, $d_y=3$, 48 systems) & fails in 40 of 48, largest gap 0.058 nats & no (we expected a pass, low confidence) \\
H2: the drive projection is the best projection & fails in 47 of 48, largest gap 0.27 nats & yes \\
H3: partitions do not beat $\SG$ ($d_x=d_y=1$, 27 systems) & beaten in 2 of 27 & yes \\
\bottomrule
\end{tabular}
\end{table}

\paragraph{Results.} The optimizer agrees with Theorem~\ref{thm:dx1} to $10^{-10}$ nats. For $d_x=2$, the best Gaussian
intervention is usually a noisy garbling, not a projection. A dense grid on the sphere of projections gave the same
best projection as the optimizer, to four decimals, in the five systems with the largest gaps. So the failure of H1 is
not a failure of the optimizer. This matters for the definition: the set of interventions in
Ref.~\cite{KolchinskyWolpert18} holds deterministic coarse-grainings only, and a larger set gives a lower value. The
median semantic efficiency is 0.25 (range 0.002 to 0.98). The value of information $\Delta V=V(I)-V(0)$ is negative in
41 of 48 systems. With $V=-H(X_\tau)$, the actual correlation usually makes the agent more spread out at $\tau$.

Non-Gaussian partitions reach $\Vact$ in 2 of 27 scalar systems, both with less information than the Gaussian value
(0.0027 against 0.0043 nats, and 0.0008 against 0.0538 nats). In both systems $|\Delta V|$ is small (0.004 and 0.001
nats). In the other 25 systems no partition of the three families reaches $\Vact$. So over all interventions, $\SG$ is
an upper bound on the semantic information, not its exact value. Part of this gap comes from the viability: a
non-Gaussian $X_\tau$ has a lower entropy than a Gaussian with the same variance.

\section{Observed semantic information}
\label{sec:dynamic}

\paragraph{Setting.} Here time is discrete, $z_{t+1}=Fz_t+w_t$, with $w_t\sim\mathcal{N}(0,W)$ and
$W=\mathrm{diag}(W_x,W_y)$, from the stationary start. The observed form of Ref.~\cite{KolchinskyWolpert18} (Sec.~V B)
uses the transfer entropy $\mathrm{TE}=\sum_{t=0}^{\tau-1}I(X_{t+1};Y_t\mid X_t)$ \cite{Schreiber00}. It applies one
coarse-graining at each step. The intervened channel is $p(x_{t+1}\mid x_t,\varphi(y_t))$, built from the actual
channel and the intervened joint law at time $t$. The response of the environment does not change. We use one
Gaussian sensor $Z_t=Ky_t+N_t$ for all $t$, with $K=\mathrm{diag}(\sqrt s)R^\top$, $\Sigma_N=\mathrm{diag}(1-s)$, a
rotation $R$ and $s\in[0,1]^{d_y}$. This family holds every Gaussian sensor up to an invertible change of $z$. The
intervened process stays Gaussian, and a recursion of its covariance gives $V$ and TE exactly. $\STE$ is the smallest
TE with $V=\Vact$, and $\SDI$ is the same with the directed information
$\mathrm{DI}=\sum_tI(Y_{0..t};X_{t+1}\mid X_{0..t})$ \cite{Massey90}.

\begin{proposition}[one step]\label{prop:F1}
Under the intervention of Ref.~\cite{KolchinskyWolpert18}, Eq.~(18), the law of $x_{t+1}$ given $x_t$ does not change.
So $p(x_1)$ does not change, and $\Delta V_{\mathrm{obs}}=\STE=0$ at $\tau=1$, for every system.
\end{proposition}
\begin{proof}
$\sum_y\hat p(x'\mid x,y)\hat p(y\mid x)=\sum_y p(x'\mid x,y)\hat p(y\mid x)$, and at $t=0$ the intervened joint is the
actual one.
\end{proof}

\begin{proposition}[DI and TE]\label{prop:F2}
$\mathrm{DI}\le\mathrm{TE}$ at each step of the intervened process. So $\SDI\le\STE$.
\end{proposition}
\begin{proof}
The intervened process is Markov in $(x_t,y_t)$, because each step uses new noise. So
$H(X_{t+1}\mid X_{0..t},Y_{0..t})=H(X_{t+1}\mid X_t,Y_t)$, and $H(X_{t+1}\mid X_{0..t})\le H(X_{t+1}\mid X_t)$.
\end{proof}

\paragraph{The link to control.} Tanaka et al.\ minimise the directed information from the plant to the controller at
a fixed quadratic cost \cite{TanakaEsfahaniMitter18}. The optimal policy is a memoryless linear Gaussian sensor, a
Kalman filter and a certainty-equivalence controller; see also Ref.~\cite{SabagTianKostinaHassibi23}. Our sensor class
is the same, but the decoder is the fixed agent. We measure the sensor as Tanaka et al.\ do,
$\mathrm{DT}=\sum_tI(Y_{0..t};Z_t\mid Z_{0..t-1})$, and we let $\DT$ be the smallest DT with $V=\Vact$. This is not their
problem: the decoder is fixed, the cost is an entropy, the horizon is finite, and the sensor does not change in time.
At the first step, given $X_0$, the chain $Y_0\to Z_0\to X_1$ gives
$I(Y_0;X_1\mid X_0)\le I(Y_0;Z_0)$. We have no proof for later steps. So we tested the conjecture
$\SDI\le\STE\le\DT$.

\paragraph{Tests.} Two groups of random systems, $F=0.85A/\rho(A)$ with Gaussian $A$, at $\tau=4$: D1 with $d_x=1$,
$d_y=2$ (40 systems) and D2 with $d_x=2$, $d_y=2$ (30 systems). One system of D1 was skipped ($|\Delta
V_{\mathrm{obs}}|<10^{-3}$). Table~\ref{tab:b3} gives the rules.

\begin{table}[t]
\centering
\caption{Observed semantic information: pre-registered rules (69 systems).}
\label{tab:b3}
\begin{tabular}{p{6.0cm}p{5.6cm}p{2.4cm}}
\toprule
Rule & Result & As expected? \\
\midrule
Sanity: $\tau=1$ keeps $V$ (Prop.~\ref{prop:F1}) & pass, largest gap $8.9\cdot10^{-16}$ & yes \\
Sanity: $\mathrm{TE}-\mathrm{DI}\ge0$ (Prop.~\ref{prop:F2}) & pass, smallest $5.7\cdot10^{-7}$ & yes \\
H1: projections reach $\STE$ & fails in 31 of 69 & yes (fails) \\
H2: the actual sensor is not optimal & holds in all 39 of D1, fails in 10 of 30 of D2 & yes, in both groups \\
H3: $\mathrm{TE}\le\mathrm{DT}$ at each noisy sensor & holds at all 1508 sensors, smallest gap $+0.030$ & yes \\
H4: $\DT$ attained and equal to $\STE$ & fails in all 69 ($\DT$ not attained in 11) & yes (fails) \\
\bottomrule
\end{tabular}
\end{table}

\paragraph{Results.} The link to control is a bound, not an equality: $\SDI\le\STE\le\DT$. The first inequality is
Proposition~\ref{prop:F2}. The second holds at every tested sensor, and it is proved only for the first step. Where
$\DT$ is attained, the median of $\STE/\DT$ is 0.15 in D1 and 0.17 in D2 (largest 0.47). So the sensor must carry about
six times more information than the agent uses. The median efficiency $\STE/\mathrm{TE}_{\mathrm{act}}$ is 0.56 in D1
and 0.68 in D2. The median of $\SDI/\STE$ is 0.84 in D1 and 0.75 in D2 (smallest 0.47), so the choice between DI and TE
matters. An agent with one coordinate never uses its actual sensor optimally (H2 in D1). In the 10 systems of D2 where
the actual sensor is optimal, $\DT$ is also not attained. We have no proof of this link.

\section{Application to a homeostat and to linear worlds}
\label{sec:apply}
Two facts limit what the stored form can say.

\begin{proposition}[scalar environment]\label{prop:F4}
Let $d_y=1$, and let $V$ be the negative entropy of one coordinate of the agent. If $\Delta V\ne0$, then only $P=1$
gives $V=\Vact$. So $\SG=\Iact$, and the efficiency is 1.
\end{proposition}
\begin{proof}
$P$ is a number $p\in[0,1]$, and the variance of the coordinate at $\tau$ is affine in $p$, so $V(p)$ is monotone.
\end{proof}

\begin{proposition}[the semantic threshold]\label{prop:F5}
In a linear Gaussian system, if some garbling changes $V$ to first order, the information/viability curve grows like
$\sqrt R$ near $R=0$. So the peak of $\mathrm{d}V/\mathrm{d}R$, the semantic threshold of Ref.~\cite{Sowinski23}, is
always at $R=0$.
\end{proposition}
\begin{proof}
$V(P)-V(0)$ is linear in $P$ to first order (Eq.~(\ref{eq:affine})), while $I_P$ is quadratic in $P$ near 0.
\end{proof}

\paragraph{A linear homeostat.} The companion study \cite{Beiro26} builds 400 linear organisms. Each has a hidden
environment $h$ with time scale $T_h\in[2,20]$, a sensor $s$, a memory $m$ and an actuator $u$ with loop gain $L$. A
viability variable $v$ feels the push of $h$, and $u$ can push it back. Its value of information
$\mathrm{VoI}=\log V_{\mathrm{intact}}-\log V_{\mathrm{scrambled}}$ uses simulated exit times. We computed the value of
information of Ref.~\cite{KolchinskyWolpert18}, $\DVKW=V(I)-V(0)$, with $V=-H(v_\tau)$, $\tau=1$ and the agent
$A=(s,m,u,v)$. We rebuilt the data of Ref.~\cite{Beiro26} with the same seeds, and a rule checked that the summary did
not change (395 organisms count). $\DVKW$ ranks the simulated value: $\rho(\DVKW,\mathrm{VoI})=+0.363$
[$+0.274$, $+0.449$]. But its sign agrees with the sign of VoI in only 0.471 of the organisms, while VoI is positive in
0.990. $\DVKW$ is positive only when the loop gain is large (median $L$ 1.15, against 0.46 when $\DVKW\le0$). It is
not a better proxy of VoI than the predictive information $P=I(A_t;h_{t+1})$: the difference of the two rank
correlations is $-0.045$ [$-0.173$, $+0.074$]. With $V$ on the whole agent, the rank correlation is higher, $+0.651$
[$+0.593$, $+0.709$]. Because $d_y=1$, Proposition~\ref{prop:F4} makes the stored semantic information equal to the full
information, so we do not use it as a score here.

\paragraph{Structures in linear worlds.} The companion study \cite{Beiro26} defines a structure as a subset of the
variables of a 6-dimensional Ornstein--Uhlenbeck world, and its viability $V$ as its mean exit time from a typical set.
It found that a unit-free information score $\PhiTC=P+C-0.25\,\TC$ (prediction, memory and total correlation) ranks
structures by $\log V$. We asked whether $\SG$ ranks them too (90 worlds, 5580 structures, $\tau=1$, $V=-H(X_\tau)$ of
the whole structure). The score is the Spearman correlation with $\log V$ inside each world, averaged over worlds, with
a cluster bootstrap over worlds.
\begin{itemize}
\item $\SG$ ranks structures, weakly: $\rho=+0.197$ [$+0.167$, $+0.226$].
\item $\DVKW$ has the wrong sign: $\rho=-0.161$ [$-0.205$, $-0.112$]. It is positive in only 0.311 of structures.
\item $\PhiTC$ ranks better than $\SG$: by $+0.146$ [$+0.085$, $+0.209$] over all structures, and by $+0.402$
[$+0.344$, $+0.464$] among structures of equal size.
\item In equilibrium worlds, $\SG$ ranks structures ($+0.198$) and $\PhiTC$ does not ($+0.006$). This was not a rule,
so it is a description only.
\end{itemize}

\section{Negative controls}
\label{sec:controls}
A score of the use of information must give zero, or a low value, to a system that holds information but does not use
it. Our data has two such systems. We test three scores on them.

\paragraph{Causal Leverage Density.} Bartlett \cite{Bartlett24} and Sowinski et al.\ \cite{SowinskiEtAl25} drop the
viability. They scramble part of the information of a known state and run the same dynamics. Then they compare the two laws
$p$ and $\hat p$ of the state at a later time:
\begin{equation}
  \chi = \JS(p\,\Vert\,\hat p)/\Oscr ,
\end{equation}
with the Jensen--Shannon divergence in bits and $\Oscr$ the number of scrambled bits. $\chi_1(t)$ is this value at time
$t$, and $\chi_3(t_1,t_2)$ is its time average \cite{SowinskiEtAl25}, Eqs.~(4)--(7). Neither paper computes CLD for a
model or tests it on a negative control. The papers leave the scrambled degrees of freedom, $\Oscr$ for continuous
variables and the times open. Our choices follow.

\subsection{An organism that does not act}
In the homeostat of Sec.~\ref{sec:apply} without feedback ($L=0$), the sensor and the memory still follow $h$, but the
actuator ignores them. Nothing uses the information.

\begin{proposition}[no feedback]\label{prop:F6}
In the homeostat with $L=0$, the value of information on $v$ is negative for every $\tau>0$. So with
Proposition~\ref{prop:F4}, $\SKW=\Iact$ and the efficiency is 1.
\end{proposition}
\begin{proof}
$v_\tau=\phi_{vv}v_0+\phi_{vh}h_0+\phi_{vu}u_0+(\text{terms independent of the start})$, with
$\phi_{vv}=e^{-k_v\tau}>0$ and $\phi_{vh}=\int_0^\tau e^{-k_v(\tau-s)}e^{-s/T_h}\mathrm{d}s>0$. With $L=0$, $u$ is
independent of $h$, and $\mathrm{Cov}(v_0,h_0)=\int_0^\infty e^{-k_vs}e^{-s/T_h}\mathrm{d}s>0$ for $\mathrm{var}(h)=1$.
The full scrambling sets $\mathrm{Cov}(A_0,h_0)$ to zero, so $\mathrm{Var}(v_\tau)$ falls by
$2\phi_{vv}\phi_{vh}\mathrm{Cov}(v_0,h_0)>0$, and $V=-H(v_\tau)$ rises.
\end{proof}

This uses the equality $V=\Vact$ of Ref.~\cite{KolchinskyWolpert18}, Eq.~(9). With the weaker condition $V\ge\Vact$,
$P=0$ is feasible when $\DVKW<0$, and $\SKW=0$. For CLD we used the full scrambling of the joint state $(A,h)$ at
$t=0$ and $\Oscr=I(A_0;h_0)$ in bits. JS comes from the exact Gaussian laws, by Monte Carlo with $2\cdot10^5$ samples
from each law.

\begin{table}[t]
\centering
\caption{The homeostat without feedback (395 organisms, $\tau=1$). N4 is the fraction of organisms where the intact
version scores higher, with 1000 bootstrap resamples of organisms.}
\label{tab:c1}
\begin{tabular}{p{4.6cm}p{6.4cm}p{3.0cm}}
\toprule
Rule & Result & Verdict \\
\midrule
N1: $\DVKW=0$? & $\DVKW<0$ in all, median $-0.255$ & not zero \\
N2: $\SKW=0$? & smallest 0.283 nats, efficiency 1 in all & not zero \\
N3: CLD $\chi_1=0$? & smallest 0.030 bits ($\ge60$ SE), median 0.094 & not zero \\
N4: $\DVKW$, intact higher & 0.787 [0.747, 0.823] & intact higher \\
N4: $\SKW$, intact higher & 0.433 [0.385, 0.481] & no feedback higher \\
N4: CLD $\chi_1$, intact higher & 0.600 [0.552, 0.646] & intact higher \\
\bottomrule
\end{tabular}
\end{table}

\paragraph{Results.} Table~\ref{tab:c1} gives the rules. All expectations were met. No score gives zero. $\DVKW$ is
negative, so its sign carries the control, but its size does not. It ranks the intact organism higher in only 79\% of
cases. It is also negative in 53.4\% of intact organisms. $\SKW$ is the full information in both versions, and it is
a little higher without feedback. CLD gives the same median to both versions (0.095 intact and 0.094 without feedback).
The sensor and the memory follow $h$ in both versions, and CLD sees this correlation. With the condition
$V\ge\Vact$, $\SKW$ is zero without feedback, but also in 53.4\% of the intact organisms. The simulated VoI is zero
without feedback by construction, because the actuator reads nothing.

\subsection{Passive memory}
The companion study \cite{Beiro26} built inertial units $z=(x,v)$ with noise. \emph{Rocks} sit in a double well
$U=h(x^2-1)^2$ with constant friction. They persist because they stay in a well, with no drive. \emph{Oscillators}
have a harmonic potential and an active friction that pumps energy in at low speed
\cite{SchweitzerEbelingTilch98,Romanczuk12}. Their life $\Von$ is the time to leave the cycle. With the drive off,
their life $\Voff$ is about 23 times shorter. Neither class uses information about an environment, but only the
oscillator does work to persist. That study used 7 draws of 60 rocks and 60 oscillators. There, the memory
$C=I(z_t;z_{t+\tau^\ast})$ at a long common lag rated rocks slightly above oscillators of the same life (stratified AUC
0.57 [0.51, 0.65]).

For CLD, we scramble one bit, $b=\mathrm{sign}(x)$. The start is the stationary law of the unit with $x>0$, so $b$ is
known. The intervention sets $b$ to a fair coin, which maps each walker to $(-x,-v)$ with probability $\tfrac12$. Both
classes are symmetric under this map, so the intervention does not change the dynamics, and
$\hat p=\tfrac12(p+Rp)$ with $R(x,v)=(-x,-v)$. We estimate JS from one $40\times40$ histogram, symmetric about 0, with
$2\cdot10^4$ walkers. A second group of walkers starts with $b$ already scrambled. Its true CLD is zero, so its value is
the bias of the histogram, and we subtract it. $\tau^\ast$ is the median rock life of the draw, and $\chi_3$ is the
average over 41 equal steps on $[0,\tau^\ast]$. The rule compares rocks and oscillators of the same life. It
computes AUC $=P(\text{rock}>\text{oscillator})$ inside tertiles of $\log V$, in bins with at least 5 units of each
class. Then it averages over bins, and over the 7 draws with a bootstrap over draws.

\begin{table}[t]
\centering
\caption{Rocks against oscillators (7 draws, 840 units). The rebuilt units equal the published summaries in all
draws.}
\label{tab:c2}
\begin{tabular}{p{5.4cm}p{4.6cm}p{3.8cm}}
\toprule
Rule & Result & Verdict \\
\midrule
K1 (main): stratified AUC of $\chi_3(0,\tau^\ast)$ & 0.711 [0.676, 0.748] & CLD mislabels rocks \\
K2: stratified AUC of $\chi_1(\tau^\ast)$ & 0.632 [0.576, 0.696] & CLD mislabels rocks \\
K3: $\rho(\chi_3,\log V)$ over rocks & 0.997 [0.996, 0.998] & yes \\
Gate: does CLD pass the control? & & fails \\
\bottomrule
\end{tabular}
\end{table}

\paragraph{Results.} Table~\ref{tab:c2} gives the rules. CLD rates rocks above oscillators of the same life in all 7
draws (AUC 0.646 to 0.784 per draw). Among rocks, $\chi_3$ follows $\log V$ almost exactly. So $\chi_3$ measures how
long the bit stays in memory: a rock keeps its well for a time of the order of its life, and an oscillator loses its
phase earlier. CLD mislabels rocks more strongly than memory did (0.711 against 0.57). We expected an AUC above 0.8, so
the effect is weaker than we predicted, but the verdicts do not depend on this. At $\tau^\ast$ itself, CLD is small
(median $\chi_1$ 0.002 to 0.003 bits for rocks), below the bias of the histogram (about 0.005 bits). So K2 rests on a
small difference, and K1, over the whole interval, is the main rule. At the own lag $V/2$ of each unit, CLD rates
oscillators higher (mean AUC 0.341), as memory did in the companion study. There, each unit has a different lag, so
this is not a comparison at the same life.

\section{Discussion}
\label{sec:discussion}

\paragraph{What the Gaussian case gives.} In linear Gaussian systems the stored semantic information of Kolchinsky and
Wolpert is a convex search over one matrix, with a closed form for an agent with one coordinate. The observed form is
exact through a covariance recursion. This removes the cost that earlier work named as the obstacle. Two results bear
on the definition itself. First, the class of interventions matters: noisy garblings beat deterministic
coarse-grainings in most systems with two agent coordinates, and non-Gaussian partitions can beat both. Second, one
step of the observed form never changes the viability, so the observed semantic information comes from correlations
that the agent builds over two or more steps.

\paragraph{What the controls show.} Table~\ref{tab:summary} collects the verdicts. No score gives zero for a system that
holds information without use. The value of information of Kolchinsky and Wolpert has the correct order in the
homeostat, but it is negative in most systems: correlation with the environment usually widens the agent. As a score of
persistence in linear worlds it has the wrong sign. The stored semantic information is the full mutual information
whenever the environment is one number and the viability is one coordinate. CLD gives the same value to an organism
that acts and to one that does not, and it prefers passive memory to driven persistence. These scores measure
correlation or memory. The protocol of Ref.~\cite{SowinskiFrankGhoshal24} asks the user to choose the viability and the
intervention, and our results show that these choices decide whether a score passes a negative control.

\begin{table}[t]
\centering
\caption{Negative controls: summary.}
\label{tab:summary}
\begin{tabular}{llll}
\toprule
Score & Control & Zero or low? & Verdict \\
\midrule
$\DVKW$ & homeostat without feedback & no, negative; order correct (0.787) & fails as a zero test \\
$\SKW$ & homeostat without feedback & no, equals $\Iact$ & fails \\
CLD & homeostat without feedback & no, same median as intact & fails \\
CLD & rocks against oscillators & no, rocks higher (0.711) & fails \\
\bottomrule
\end{tabular}
\end{table}

\paragraph{What passes.} In each control, a term that sees the drive or the loop separates the classes. In the
homeostat, the simulated VoI is zero without feedback, because it scrambles the input of the actuator at all times.
For the units, the flux ratio $\Von/\Voff$ separates rocks from oscillators with no overlap \cite{Beiro26}: a rock lives
the same with the drive off, and an oscillator lives much less. Both need an intervention on the mechanism, not only on
the correlations of a state. We propose these two systems, with code, as negative controls for any score of semantic
information or agency.

\paragraph{Limits.} The viability $-H(X_\tau)$ follows Ref.~\cite{KolchinskyWolpert18}, and it is not an exit time. Our
Gaussian values are upper bounds over all interventions. The bound $\STE\le\DT$ is proved only for the first step. We
fix one choice for CLD where the papers leave it open (the scrambled bit, $\Oscr$, the times and the JS functional).
Another choice can give another result. The random systems come from one construction each, with one value of $\tau$
in each rule.

\paragraph{Open problems.} A proof or a counterexample for $\mathrm{TE}\le\mathrm{DT}$ after the first step. A closed
form, or a tight bound, for $d_x\ge2$. The gap to Tanaka et al.\ with a Kalman filter as decoder. And a score of use
that passes both controls without a second simulation.

\section*{Data and code availability}
The code, the step documents with the rules, the reports and the logs are public at
\url{https://github.com/PepeDest/Beiro-Feat-Flutter-Big-ass-project}. The formulas are in
\texttt{src/vlife/semantic.py}, \texttt{semantic\_dynamic.py}, \texttt{semantic\_apply.py} and
\texttt{negative\_controls.py}. The runs are \texttt{experiments/e20} to \texttt{e26}. The test suite has 113 tests.

\section*{Deviations from the rules}
\begin{itemize}
\item B1: the pass rule did not allow for the bias of the plug-in estimator at zero information. We report the
verdict ``fail'' as written (Sec.~\ref{sec:b1}).
\item B1: a garbling matrix in the simulation script acted on raw $y$ instead of whitened $y$. We fixed it before the
run.
\item B4: one script stopped at the last print, after it wrote its results. The second run has no random step and gave
the same numbers.
\item C2: the first run stopped when the computer shut down. A resume option reused the saved units, checked them again
against the summaries, and saved the CLD results after each draw. The rules did not change.
\end{itemize}

\section*{Use of generative AI}
The code, the analyses and a first draft of this text were produced with substantial help from Claude (Anthropic), a
generative AI model, under the direction of the author. The author takes full responsibility for all contents of this
paper.

\bibliographystyle{unsrtnat}
\bibliography{refs}

\end{document}
